\section{Combining Setups and Context}

\subsection{The sources describe trades as assembled, not selected}

One of the strongest messages in the transcript material is that a good trade is built like a constructor set: multiple pieces reinforce each other. That idea is essential for both discretionary trading and future research design. It means the framework should not ask only \emph{``Did setup X appear?''} It should ask:

\begin{itemize}[leftmargin=1.5em]
\item What is the primary setup family?
\item What evidence supports it from other modalities?
\item What evidence contradicts it?
\item How much of the evidence is structural versus ephemeral?
\item How fragile is the idea if one ingredient disappears?
\end{itemize}

\subsection{When a bounce becomes a breakout}

A bounce setup can mutate into a breakout setup when the same defense is tested too often or is visibly degrading. The transformation usually follows this sequence:

\begin{enumerate}[leftmargin=1.5em]
\item a defended level first produces good rejections;
\item the market returns to it repeatedly;
\item each test consumes more liquidity or shortens the distance between tests;
\item tape aggression remains persistent;
\item the defending side looks less able to repel the market cleanly;
\item once breached, the level flips from barrier to fuel source.
\end{enumerate}

For research, this implies that ``bounce'' and ``breakout'' are not mutually exclusive static labels. They are states in a local evolution process.

\subsection{When a breakout becomes a false breakout}

The failure sequence is the mirror image:

\begin{enumerate}[leftmargin=1.5em]
\item an obvious level is breached;
\item breakout participants and stops activate;
\item post-break continuation is weak or short-lived;
\item the market returns through the level;
\item trapped traders now provide reversal fuel.
\end{enumerate}

This is why the first seconds after a breach matter so much in the source material. The trader is not only asking whether price crossed the level, but whether the market \emph{accepted} the new price area.

\subsection{How order book, tape, and chart evidence should be combined}

The combined logic implied by the materials can be summarized in Table~\ref{tab:evidence_roles}.

\begin{longtable}{L{0.22\textwidth}L{0.7\textwidth}}
\caption{Role of each evidence source in the trading logic}\label{tab:evidence_roles}\\
\toprule
Evidence source & Primary role \\
\midrule
\endfirsthead
\toprule
Evidence source & Primary role \\
\midrule
\endhead
DOM / book & Shows displayed intentions, queue sizes, spread condition, visible defense, robot behavior, and microstructure asymmetry near the current price. \\
Tape / prints & Shows who is acting aggressively now, how intense the aggression is, and whether aggression actually achieves price movement. \\
Chart / levels & Provides memory and geometry: prior highs/lows, round numbers, compression, trend phase, and locally obvious liquidity pools. \\
Clusters / POC & Confirm whether important volume historically transacted where the current setup claims significance. \\
Cross-market context & Indicates whether another instrument or market is currently leading, confirming, or contradicting the setup. \\
Risk state & Determines whether the trader is still allowed to express the idea, even if the idea itself is valid. \\
\bottomrule
\end{longtable}

\subsection{Confidence as a weighted judgment, not a binary tag}

The repository should eventually represent setup quality as a weighted confidence object, not a one-bit flag. The source materials imply at least the following dimensions:

\begin{itemize}[leftmargin=1.5em]
\item \textbf{Structural quality:} level significance, round number, chart alignment, historical footprint support.
\item \textbf{Microstructure quality:} queue behavior, spread, aggression, replenishment, and whether price is responding properly.
\item \textbf{Context quality:} active instrument, relevant time of day, event sensitivity, leader-follower confirmation.
\item \textbf{Risk quality:} whether the stop is obvious and compact, whether slippage is manageable, and whether the trader has remaining daily risk budget.
\end{itemize}

\subsection{Day-specific adaptation}

The transcript on trade review makes an especially important point: what worked yesterday may not work today, and what failed for a month may come back. This is more than trader psychology; it is a research direction. The framework should track which setup families are currently being rewarded for each instrument on the current day.

A practical interpretation is:

\begin{enumerate}[leftmargin=1.5em]
\item Observe the opening phase and early session.
\item Log all high-quality setup candidates, even if not traded.
\item Score their near-term outcomes using simple rules.
\item Maintain a live same-day setup quality profile.
\item Use that profile as a filter or confidence modifier for later signals.
\end{enumerate}

This is one of the most valuable bridges from discretionary practice to systematic research in the entire knowledge base.
